\section{AI 应用方法论：模型能力、场景补充与外溢}

上一章讲 Monica，本章抽象出肖弘的 AI 应用分类。他把火过的 AI 应用当作少量数据点，总结三类机会：第一，主场景补充，例如在搜索、写作、浏览器、办公等已有场景里增强体验；第二，模型能力带来的新玩法，尤其在图像和视频生成里容易爆发；第三，模型能力在特定领域外溢，解决了该领域核心问题，只缺一个好容器承接，比如 Cursor 承接 coding 能力。

\begin{figure}[H]
\centering
\includegraphics[width=0.92\textwidth]{ai-app-taxonomy.png}
\caption{AI 应用分类方法：主场景补充、模型能力变化、特定领域外溢。自制概念图，依据 02:02:14--02:09:32 对谈内容整理。}
\end{figure}

\begin{knowledgebox}{读图：找机会不是追热点，而是找能力和场景的缺口}
主场景补充看现有工作流哪里缺 AI；能力变化看新模型打开的新玩法；领域外溢看通用模型是否已经解决某个垂直领域最难的部分，而产品公司只需要把边缘工作流补齐。
\end{knowledgebox}

\subsection{Cursor 例子：核心能力外溢与容器}

本节展开 Cursor。肖弘说，他曾用 AI 对话生成代码，再手动复制粘贴、新建文件、对比差异；后来意识到 Cursor 做对的是把这些边缘动作工程化。写代码的核心能力来自 foundation model，产品公司不必重新解决；但它可以提供编辑器容器、文件操作、差异对比、上下文管理和顺滑交互。这就是“模型能力外溢后，容器承接”的典型例子。

\begin{importantbox}{AI 应用机会公式}
如果通用模型已经解决某领域的核心能力，但用户仍然要手工完成大量边缘步骤，那么应用公司可以做容器：把模型能力接入真实工作流，把复制、粘贴、对比、调用工具、保存结果这些摩擦消掉。
\end{importantbox}

\subsection{大模型原厂 vs 应用公司}

本节区分两类公司。模型原厂能决定能力什么时候发生，因此有更强控制力；应用公司无法决定模型能力边界，只能预测、等待、快速接入。但应用公司也有优势：可以使用全世界模型，专注特定场景、用户体验、品牌、增长和反馈。肖弘把这看作“贸工技”路线：先做用户和产品，再逐步补工程和特定模型能力。

\begin{figure}[H]
\centering
\includegraphics[width=0.92\textwidth]{model-company-vs-app-company.png}
\caption{大模型原厂 vs 应用公司：原厂做通用能力，应用公司做场景、容器和反馈。自制概念图，依据 02:21:01--02:23:51 对谈内容整理。}
\end{figure}

\begin{knowledgebox}{读图：控制能力发生 vs 承接能力外溢}
模型原厂能推进底座能力，甚至定义能力何时出现；应用公司更像站在能力浪潮边上，提前摆好产品容器。两者不是简单上下游，而是控制权、速度、用户理解和场景深度的不同组合。
\end{knowledgebox}

\subsection{DeepSeek 的启发}

本节看 DeepSeek 对中国 AI 应用生态的冲击。肖弘说 DeepSeek “最佛却取得最好结果”，给应用创业者精神鼓励：少一点表演，多一点真实能力。DeepSeek 爆火也让他重新思考产品传播：OpenAI o1 有强能力但产品表达没抓住大众心智，而 DeepSeek 用更低姿态和更直接可用的方式冲击全球用户。

\begin{figure}[H]
\centering
\includegraphics[width=0.90\textwidth]{deepseek-product-lesson.png}
\caption{DeepSeek 产品启发：最佛的团队取得最好结果，给应用创业者精神鼓励。自制概念图，依据 02:23:51--02:29:00 对谈内容整理。}
\end{figure}

\begin{knowledgebox}{读图：能力与姿态一起改变传播}
DeepSeek 一侧强调能力、开源和低姿态；应用创业一侧强调真实价值。这里的 lesson 不是所有公司都要佛系，而是产品传播不能只靠营销姿态，底层能力和用户感知必须站得住。
\end{knowledgebox}

\subsection{DeepSeek vs o1：能力强，还要被用户感知}

本节补充第二次访谈里的产品视角。肖弘谈 DeepSeek 时，并不是在做模型评测，而是在分析一个模型能力事件如何变成用户心智事件。OpenAI o1 在推理能力上有重要突破，但它的产品表达更偏“模型版本”和“专业能力”；DeepSeek 则以更低门槛、更开放、更容易传播的方式进入普通用户视野。对应用创业者来说，这提醒他们：技术能力本身不自动等于产品爆发，能力必须被包装成用户能直接理解、直接试用、直接传播的体验。

\begin{importantbox}{产品表达的三层}
第一层是能力真实存在；第二层是用户能低成本触达和验证；第三层是用户能用一句话讲给别人。DeepSeek 之所以给应用创业者精神鼓励，是因为它把能力、开放姿态和传播路径同时打通。
\end{importantbox}

\noindent\begin{tabularx}{\textwidth}{p{0.20\textwidth}p{0.34\textwidth}X}
\toprule
维度 & o1 式表达 & DeepSeek 式表达 \\
\midrule
用户感知 & 更像高阶推理模型版本 & 更像人人可试的能力事件。 \\
传播方式 & 专业圈层和 benchmark 讨论更强 & 开源、低成本、中文社区传播更强。 \\
应用启发 & 能力强但产品包装不一定充分 & 能力和分发结合，给应用层带来信心。 \\
\bottomrule
\end{tabularx}

\subsection{本章小结}

AI 应用方法论不是“找个模型套壳”。它要求判断模型能力的下一步、找到用户主场景、补齐容器摩擦、持续理解用户，并接受产品形态会随着模型能力不断重写。

